% ---------------------------------------------------------------------------
% French edition: layout adjustments, loaded after header.tex.
% pandoc declares the language with \babelprovide[main,import]{french}, which
% takes the French names (Partie, Chapitre, Annexe, Table des matières) and the
% hyphenation patterns from babel's locale files without loading babel-french's
% typographic extensions. The French punctuation spacing therefore comes from
% the Markdown sources themselves (narrow no-break spaces, which the Lua filter
% turns into \,), and only the figure caption needs a French colon.
% ---------------------------------------------------------------------------
\DeclareCaptionLabelSeparator{frcolon}{\,: }
\captionsetup{labelsep=frcolon}
% babel's French locale spells the part label in lower case ("partie I"); the
% part pages of the book carry it capitalised, like "Chapitre" and "Annexe".
\AtBeginDocument{\renewcommand{\partname}{Partie}}
